% GENERATED FILE. Edit canonical lessons, then run npm run generate:books.
% canonical-source-hash: fnv1a64:f948b089b7789da6
% canonical-lessons: MR-C434-gairsoy, MR-C434-cukamuk, MR-C434-phajiti, MR-C434-dhavpal-karne, MR-C434-marga-kadhne

\chapter{Mix-ups and Fiascos}
\label{ch:mr-mix-ups-and-fiascos}
\hlchaptermodality{434}

\begin{chapteropening}
\textbf{By the end of this chapter:} \emph{I can talk about inconvenience, missed meetings and fiascos, and finding a way through.}

Complete the last lesson of chapter 434: I can talk about inconvenience, missed meetings and fiascos, and finding a way through.
\end{chapteropening}

\section[\texorpdfstring{\mr{गैरसोय} (gairsoy)}{गैरसोय (gairsoy)}]{\mr{गैरसोय} (gairsoy) — inconvenience}

\label{lesson:MR-C434-gairsoy}



\begin{quote}
Before the new one: say the Marathi for to get dented, then the Marathi for to stop working.
\end{quote}

\subsection*{You'll want to know: \mr{गैरसोय}}
\textbf{\mr{गैरसोय}} — \emph{gairsoy} — \textquotedblleft{}inconvenience\textquotedblright{}.

\begin{quote}
\mr{पाणी} \mr{नसल्यामुळे} \mr{सगळ्यांची} \mr{गैरसोय} \mr{झाली}.
\end{quote}

\begin{quote}
\emph{pāṇī nasalyāmuḷe sagaḷyāñcī gairsoy jhālī.}
\end{quote}

\begin{quote}
Everyone was put out because there was no water.
\end{quote}

Feminine.

\mr{गैरसोय} \mr{होणे}, to be inconvenienced: the person takes -\mr{ची}.

\begin{cousinweb}
\textbf{\mr{गैरसोय}} (\emph{gairsoy}) is built from \emph{gair-}, \textquotedblleft{}un-, other than\textquotedblright{}, a loan usually traced to Arabic \emph{ġair}, and \textbf{\mr{सोय}}, \textquotedblleft{}convenience\textquotedblright{}.
\end{cousinweb}

\begin{grammarlens}[title={What you've built}]
One more word for getting around.
\end{grammarlens}

\subsection*{Your turn}
\begin{itemize}
\raggedright
  \item \emph{Say it:} \emph{gairsoy}
  \item \emph{Say it:} \emph{gairsoy}, once more
  \item \emph{Say it:} \emph{pāṇī nasalyāmuḷe sagaḷyāñcī gairsoy jhālī.} — Everyone was put out because there was no water.
  \item \emph{Recall:} say the Marathi for to get dented, then the Marathi for to stop working, then say \emph{gairsoy} again
\end{itemize}

\begin{tcolorbox}[breakable,colback=teal!4,colframe=teal!35!black,title={Before you move on}]
What does \mr{गैरसोय} mean? (\textquotedblleft{}Inconvenience\textquotedblright{}.) Say it once more.
\end{tcolorbox}
\section[\texorpdfstring{\mr{चुकामूक} (cukāmūk)}{चुकामूक (cukāmūk)}]{\mr{चुकामूक} (cukāmūk) — missing each other}

\label{lesson:MR-C434-cukamuk}



\begin{quote}
Before the new one: say the Marathi for to stop working, then the Marathi for inconvenience.
\end{quote}

\subsection*{You'll want to know: \mr{चुकामूक}}
\textbf{\mr{चुकामूक}} — \emph{cukāmūk} — \textquotedblleft{}missing each other\textquotedblright{}.

\begin{quote}
\mr{स्टेशनवर} \mr{आमची} \mr{चुकामूक} \mr{झाली}.
\end{quote}

\begin{quote}
\emph{sṭeśanvar āmcī cukāmūk jhālī.}
\end{quote}

\begin{quote}
We missed each other at the station.
\end{quote}

Feminine.

\mr{चुकामूक} \mr{होणे}, two people fail to meet.

\begin{cousinweb}
\textbf{\mr{चुकामूक}} (\emph{cukāmūk}) is an echo compound built on \textbf{\mr{चुकणे}}, \textquotedblleft{}to miss\textquotedblright{}, whose origin is not settled.
\end{cousinweb}

\begin{grammarlens}[title={What you've built}]
One more word for getting around.
\end{grammarlens}

\subsection*{Your turn}
\begin{itemize}
\raggedright
  \item \emph{Say it:} \emph{cukāmūk}
  \item \emph{Say it:} \emph{cukāmūk}, once more
  \item \emph{Say it:} \emph{sṭeśanvar āmcī cukāmūk jhālī.} — We missed each other at the station.
  \item \emph{Recall:} say the Marathi for to stop working, then the Marathi for inconvenience, then say \emph{cukāmūk} again
\end{itemize}

\begin{tcolorbox}[breakable,colback=teal!4,colframe=teal!35!black,title={Before you move on}]
What does \mr{चुकामूक} mean? (\textquotedblleft{}Missing each other\textquotedblright{}.) Say it once more.
\end{tcolorbox}
\section[\texorpdfstring{\mr{फजिती} (phajitī)}{फजिती (phajitī)}]{\mr{फजिती} (phajitī) — a fiasco; an embarrassment}

\label{lesson:MR-C434-phajiti}



\begin{quote}
Before the new one: say the Marathi for inconvenience, then the Marathi for missing each other.
\end{quote}

\subsection*{You'll want to know: \mr{फजिती}}
\textbf{\mr{फजिती}} — \emph{phajitī} — \textquotedblleft{}a fiasco; an embarrassment\textquotedblright{}.

\begin{quote}
\mr{पावसात} \mr{छत्री} \mr{उलटली} \mr{आणि} \mr{माझी} \mr{फजिती} \mr{झाली}.
\end{quote}

\begin{quote}
\emph{pāvsāt chatrī ulaṭlī āṇi mājhī phajitī jhālī.}
\end{quote}

\begin{quote}
In the rain my umbrella turned inside out and I was embarrassed.
\end{quote}

Feminine.

\mr{फजिती} \mr{होणे}, to be made to look foolish; the person takes -\mr{ची}.

\begin{cousinweb}
\textbf{\mr{फजिती}} (\emph{phajitī}) is a loan through Persian, usually traced to Arabic \emph{faḍīḥat}, \textquotedblleft{}disgrace\textquotedblright{}.
\end{cousinweb}

\begin{grammarlens}[title={What you've built}]
One more word for getting around.
\end{grammarlens}

\subsection*{Your turn}
\begin{itemize}
\raggedright
  \item \emph{Say it:} \emph{phajitī}
  \item \emph{Say it:} \emph{phajitī}, once more
  \item \emph{Say it:} \emph{pāvsāt chatrī ulaṭlī āṇi mājhī phajitī jhālī.} — In the rain my umbrella turned inside out and I was embarrassed.
  \item \emph{Recall:} say the Marathi for inconvenience, then the Marathi for missing each other, then say \emph{phajitī} again
\end{itemize}

\begin{tcolorbox}[breakable,colback=teal!4,colframe=teal!35!black,title={Before you move on}]
What does \mr{फजिती} mean? (\textquotedblleft{}A fiasco; an embarrassment\textquotedblright{}.) Say it once more.
\end{tcolorbox}
\section[\texorpdfstring{\mr{धावपळ} \mr{करणे} (dhāvpaḷ karṇe)}{धावपळ करणे (dhāvpaḷ karṇe)}]{\mr{धावपळ} \mr{करणे} (dhāvpaḷ karṇe) — to rush about}

\label{lesson:MR-C434-dhavpal-karne}



\begin{quote}
Before the new one: say the Marathi for missing each other, then the Marathi for a fiasco.
\end{quote}

\subsection*{You'll want to know: \mr{धावपळ} \mr{करणे}}
\textbf{\mr{धावपळ} \mr{करणे}} — \emph{dhāvpaḷ karṇe} — \textquotedblleft{}to rush about\textquotedblright{}.

\begin{quote}
\mr{गाडी} \mr{पकडण्यासाठी} \mr{आम्ही} \mr{खूप} \mr{धावपळ} \mr{केली}.
\end{quote}

\begin{quote}
\emph{gāḍī pakḍaṇyāsāṭhī āmhī khūp dhāvpaḷ kelī.}
\end{quote}

\begin{quote}
We rushed about a lot to catch the train.
\end{quote}

Past, with -\mr{ने} (he …): \mr{त्याने} \mr{धावपळ} \mr{केली}, \emph{tyāne dhāvpaḷ kelī}.

\mr{धावपळ} is feminine, so the past is \mr{धावपळ} \mr{केली}.

\begin{cousinweb}
\textbf{\mr{धावपळ} \mr{करणे}} (\emph{dhāvpaḷ karṇe}) is built from \textbf{\mr{धाव}}, \textquotedblleft{}a run\textquotedblright{} (CDIAL 6800), and \emph{paḷ}, the stem of \textbf{\mr{पळणे}}, \textquotedblleft{}to run off\textquotedblright{}, with \textbf{\mr{करणे}}.
\end{cousinweb}

\begin{grammarlens}[title={What you've built}]
One more word for getting around.
\end{grammarlens}

\subsection*{Your turn}
\begin{itemize}
\raggedright
  \item \emph{Say it:} \emph{dhāvpaḷ karṇe}
  \item \emph{Say it:} \emph{dhāvpaḷ karṇe}, once more
  \item \emph{Say it:} \emph{gāḍī pakḍaṇyāsāṭhī āmhī khūp dhāvpaḷ kelī.} — We rushed about a lot to catch the train.
  \item \emph{Recall:} say the Marathi for missing each other, then the Marathi for a fiasco, then say \emph{dhāvpaḷ karṇe} again
\end{itemize}

\begin{tcolorbox}[breakable,colback=teal!4,colframe=teal!35!black,title={Before you move on}]
What does \mr{धावपळ} \mr{करणे} mean? (\textquotedblleft{}To rush about\textquotedblright{}.) Say it once more.
\end{tcolorbox}
\section[\texorpdfstring{\mr{मार्ग} \mr{काढणे} (mārga kāḍhṇe)}{मार्ग काढणे (mārga kāḍhṇe)}]{\mr{मार्ग} \mr{काढणे} (mārga kāḍhṇe) — to find a way; to get round a problem}

\label{lesson:MR-C434-marga-kadhne}



\begin{quote}
Before the new one: say the Marathi for a fiasco, then the Marathi for to rush about.
\end{quote}

\subsection*{You'll want to know: \mr{मार्ग} \mr{काढणे}}
\textbf{\mr{मार्ग} \mr{काढणे}} — \emph{mārga kāḍhṇe} — \textquotedblleft{}to find a way; to get round a problem\textquotedblright{}.

\begin{quote}
\mr{पुलावर} \mr{कोंडी} \mr{होती}, \mr{पण} \mr{चालकाने} \mr{मार्ग} \mr{काढला}.
\end{quote}

\begin{quote}
\emph{pulāvar koṇḍī hotī, paṇ cālakāne mārga kāḍhlā.}
\end{quote}

\begin{quote}
The bridge was jammed, but the driver found a way through.
\end{quote}

Past, with -\mr{ने} (he …): \mr{त्याने} \mr{मार्ग} \mr{काढला}, \emph{tyāne mārga kāḍhlā}.

Also to solve a problem: \mr{अडचणीतून} \mr{मार्ग} \mr{काढणे}.

\begin{cousinweb}
\textbf{\mr{मार्ग} \mr{काढणे}} (\emph{mārga kāḍhṇe}) joins \textbf{\mr{मार्ग}}, \textquotedblleft{}a way\textquotedblright{}, a tatsama from Sanskrit \emph{mārga} (Monier-Williams), to \textbf{\mr{काढणे}}, \textquotedblleft{}to take out\textquotedblright{}.
\end{cousinweb}

\begin{grammarlens}[title={What you've built}]
One more word for getting around.
\end{grammarlens}

\subsection*{Your turn}
\begin{itemize}
\raggedright
  \item \emph{Say it:} \emph{mārga kāḍhṇe}
  \item \emph{Say it:} \emph{mārga kāḍhṇe}, once more
  \item \emph{Say it:} \emph{pulāvar koṇḍī hotī, paṇ cālakāne mārga kāḍhlā.} — The bridge was jammed, but the driver found a way through.
  \item \emph{Recall:} say the Marathi for a fiasco, then the Marathi for to rush about, then say \emph{mārga kāḍhṇe} again
\end{itemize}

\begin{tcolorbox}[breakable,colback=teal!4,colframe=teal!35!black,title={Before you move on}]
What does \mr{मार्ग} \mr{काढणे} mean? (\textquotedblleft{}To find a way; to get round a problem\textquotedblright{}.) Say it once more.
\end{tcolorbox}
